% !TeX root = ../se200_identity_regimes.tex

% =============================================================================
\section{Pairwise Separation}
\label{app:sep}
% =============================================================================

The nine regimes yield $\binom{9}{2}=36$ unordered pairs.
Each pair is separated either by carrier difference or, for same-carrier split pairs,
by the canonical operation that forces the split.
Only three pairs share a reference kind and differ by identity basis:
$\LOC/\OBJ$, $\SCOPEE/\SCOPES$, and $\RULEC/\RULES$.
All other pairs are separated by carrier difference.

\begin{center}
\small
\setlength{\tabcolsep}{2pt}
\begin{tabular}{@{}l ccccccccc@{}}
\toprule
 & $\OBL$ & $\OCC$ & $\REC$ & $\LOC$ & $\OBJ$ & $\SCOPEE$ & $\SCOPES$ & $\RULEC$ & $\RULES$\\
\midrule
$\OBL$   & -- & $\bullet$ & $\bullet$ & $\bullet$ & $\bullet$ & $\bullet$ & $\bullet$ & $\bullet$ & $\bullet$\\
$\OCC$   & $\bullet$ & -- & $\bullet$ & $\bullet$ & $\bullet$ & $\bullet$ & $\bullet$ & $\bullet$ & $\bullet$\\
$\REC$   & $\bullet$ & $\bullet$ & -- & $\bullet$ & $\bullet$ & $\bullet$ & $\bullet$ & $\bullet$ & $\bullet$\\
$\LOC$  & $\bullet$ & $\bullet$ & $\bullet$ & -- & BF & $\bullet$ & $\bullet$ & $\bullet$ & $\bullet$\\
$\OBJ$  & $\bullet$ & $\bullet$ & $\bullet$ & BF & -- & $\bullet$ & $\bullet$ & $\bullet$ & $\bullet$\\
$\SCOPEE$  & $\bullet$ & $\bullet$ & $\bullet$ & $\bullet$ & $\bullet$ & -- & AD & $\bullet$ & $\bullet$\\
$\SCOPES$  & $\bullet$ & $\bullet$ & $\bullet$ & $\bullet$ & $\bullet$ & AD & -- & $\bullet$ & $\bullet$\\
$\RULEC$  & $\bullet$ & $\bullet$ & $\bullet$ & $\bullet$ & $\bullet$ & $\bullet$ & $\bullet$ & -- & RF\\
$\RULES$  & $\bullet$ & $\bullet$ & $\bullet$ & $\bullet$ & $\bullet$ & $\bullet$ & $\bullet$ & RF & --\\
\bottomrule
\end{tabular}
\end{center}

Here $\bullet$ indicates separation by carrier difference.
The entries BF, AD, and RF indicate same-kind split pairs separated by the corresponding canonical operation.

The $\LOC/\OBJ$ pair is separated by branch/fork.
A fork over a fixed locus may preserve locus identity while breaking object identity.
Thus the two regimes induce distinct identity relations.

The $\SCOPEE/\SCOPES$ pair is separated by aggregation/decomposition.
A decomposition may preserve the covered cases while changing the internal scope structure.
Thus the two regimes induce distinct identity relations.

The $\RULEC/\RULES$ pair is separated by refinement.
A refinement may preserve rule content while changing textual or internal rule structure.
Thus the two regimes induce distinct identity relations.

Therefore every pair of regimes is separated either by carrier difference or by a transformation witness.
No pair collapses.
